\documentclass[11pt]{article}
\usepackage[a4paper,margin=23mm]{geometry}
\usepackage{amsmath,amssymb,amsthm,mathtools,bm,booktabs,microtype,xurl,hyperref}
\hypersetup{colorlinks=true,linkcolor=blue,urlcolor=blue,citecolor=blue,
pdftitle={P54 source-preserving finite matching}}
\setlength{\emergencystretch}{2em}
\newcommand{\tr}{\operatorname{tr}}
\newcommand{\STr}{\operatorname{STr}}
\newcommand{\diag}{\operatorname{diag}}
\newcommand{\MS}{\overline{\mathrm{MS}}}
\newcommand{\dd}{\mathrm d}
\newcommand{\ellight}{\varphi}
\newtheorem{proposition}{Proposition}
\title{Route-F P54: Source-Preserving One-Loop Matching\\
and the Upper Heavy-Vector Yukawa/Kinetic Contribution}
\author{P54PQ-v2 four-dimensional mainline}
\date{11 September 2026}
\begin{document}
\maketitle
\begin{abstract}
We derive the full first-order transport law for the kernel, interaction
source and induced contact response $(K,J,C)$. We calculate a previously
missing upper-PS heavy-vector contribution using the actual Spin(10)
generators and cached same-action stationary Hessian. Dimensional
numerator terms supply finite constants which must not be dropped before
minimal subtraction. Scalar kinetic matching generates conjugate
bidoublet Yukawa coefficients as well as the original PS tensors.
Different-colour diquark probes reveal upper scalar-exchange contact
terms that the earlier same-colour probes could not see.
The one-loop algebra is complete; the diagram content is explicitly a
subset. The complete finite matching, fermion-inclusive physical light
state and constrained flavor fit are \emph{not} closed by this work.
\end{abstract}

\section{Fixed prescription and exact scope}
The action and field content remain P54PQ-v2:
$3\,16_F+54_{H,\mathbb R}+126_{H,\mathbb C}+10_{H,\mathbb C}
+1_{H,\mathbb C}$. The only fundamental family spurions are the two
complex symmetric matrices $h_{\rm raw},f_{\rm raw}$. No extra spatial
dimension, geometry-dependent particle mass, lattice scan, or new fit
parameter is introduced.

This calculation uses dimensional regularization (DR, not dimensional
\emph{reduction}), $d=4-2\epsilon$, $\MS$ subtraction, background-field
Landau gauge, and the existing fixed-VEV tadpole prescription. Open
fermion lines need no anomalous $\gamma_5$ trace here. The upper stationary
PS background, in the previous reference units, is
\begin{equation}
 (\omega_U,\sigma_U,v_{s,U})=
 (1.00081030592151,0,0.254277697083708).
\end{equation}
It is not the lower broken background $(1,0.1265,0.25)$. Its active
scalar parents need not have positive curvature; they are retained fields
of a Wilson matching calculation, not integrated Gaussian instabilities.
Only the 55 positive physical upper scalars are integrated at the upper
scalar step. There are 24 upper eaten directions, 248 real active parent
coordinates, and one physical PQ direction.

We read the existing content-addressed Hessian rather than recompute it.
The historical gauge benchmark gives
\begin{equation}
 g_U=0.562602889985323,\qquad
 \mu_U=0.563058770438537,\qquad M_V^2=0.317035178967757.
\end{equation}
All 24 coset vectors share this mass. These numbers do not recertify the
historical dimensionful unification scales or phenomenological viability.

The computed graphs are the heavy-vector correction to the active
Yukawa vertex, the vector contribution to the fermion kinetic metric,
and the vector--scalar contribution to the active scalar kinetic metric.
Eliminated scalar masses in the last graph are kept exactly. Retained
scalar masses are expanded to leading order in the hard region. Their
largest $|m^2|/M_V^2$ is $0.328605$, so an accurate numerical fit requires
a controlled remainder or higher matching orders in this expansion.
We do not substitute the previously computed \emph{broken-background}
scalar-cubic metric into this upper-PS package.

\section{Why the source-dependent effective action is required}
Let $I$ denote a declared collection of external current slots. Choose
the Euclidean Gaussian sign convention
\begin{equation}\label{eq:gaussian}
 S_2[\phi,I]=\frac12\phi^TK\phi-\phi^TJI-\frac12 I^TCI.
\end{equation}
Then integrating $\phi$ gives the response
\begin{equation}\label{eq:response}
 F=C+J^TK^{-1}J,\qquad S_{\rm ind}=-\frac12 I^TFI.
\end{equation}
The transpose is intentional: $\phi,I$ are real-field/current coordinates
analytically continued in momentum. It is not a Hermitian conjugate at
complex momentum. Matrix spectral positivity concerns the appropriate
physical real-energy representation, not arbitrary complex $K$.

\subsection{Exact elimination and its first derivative}
Partition retained and eliminated coordinates as
\begin{equation}
 K=\begin{pmatrix}A&B\\B^T&D\end{pmatrix},\quad
 J=\binom{J_r}{J_h},\quad R=D^{-1}.
\end{equation}
Completing the square gives
\begin{align}
 K_r'&=A-BRB^T,\label{eq:treeK}\\
 J_r'&=J_r-BRJ_h,\label{eq:treeJ}\\
 C'&=C+J_h^TRJ_h.\label{eq:treeC}
\end{align}
Write every input as $X=X_0+\ell\,\delta X+O(\ell^2)$, with a formal
loop-counting parameter $\ell$. The displayed numerical $\delta X$
already contains its $1/(16\pi^2)$ coefficient. Since
$\delta R=-R\delta D R$, the complete linearized map is
\begin{align}
 \delta K_r'={}&\delta A-\delta B RB^T-BR\delta B^T
                   +BR\delta D RB^T,\label{eq:jetK}\\
 \delta J_r'={}&\delta J_r-\delta B RJ_h-BR\delta J_h
                   +BR\delta D RJ_h,\label{eq:jetJ}\\
 \delta C'={}&\delta C+\delta J_h^TRJ_h+J_h^TR\delta J_h
                   -J_h^TR\delta D RJ_h.\label{eq:jetC}
\end{align}
All unvaried quantities here are at tree order. The response derivative is
\begin{equation}\label{eq:deltaF}
 \delta F=\delta C+\delta J^TK_0^{-1}J_0+J_0^TK_0^{-1}\delta J
 -J_0^TK_0^{-1}\delta K K_0^{-1}J_0.
\end{equation}

\begin{proposition}[Order-independent one-loop transport]
If all eliminated blocks along two nested elimination paths are
invertible, and both paths start from the same source-dependent action
and expansion, their final $(K,J,C)$ and their first variations agree.
\end{proposition}
\begin{proof}
The tree maps are the same completion of the square in the joint
eliminated variables. The resulting rational matrix functions agree
on their common invertibility domain. Differentiating that identity
gives Eqs.~\eqref{eq:jetK}--\eqref{eq:jetC}; the chain rule proves the
statement. No assumption of commuting mass blocks is used.
\end{proof}
This is an algebraic theorem, not a claim that separately truncated
local EFTs automatically have identical Wilson coefficients. Local
expansion, IR subtraction, operator reduction, and scale evolution must
also be performed consistently. Singular blocks cannot be repaired by an
unannounced pseudoinverse.

\subsection{Where the missing finite diagrams enter}
For classical heavy solution $\Phi_h^{\rm cl}[\ellight,I]$, the formal
one-loop matching functional has the structure
\begin{equation}\label{eq:hard}
 \Delta\Gamma^{(1)}[\ellight,I]=
 \left[\frac12\STr\log\mathcal M_{\rm UV}[\ellight,\Phi_h^{\rm cl},I]
 -\frac12\STr\log\mathcal M_{{\rm EFT},0}[\ellight,I]
 +\Gamma_{\rm CT}\right]_{\rm local,ren}.
\end{equation}
Here $\ellight$ denotes retained fields, not loop counting;
$\mathcal M$ is the full gauge-fixed fluctuation operator. The supertrace
includes the correct real-boson, fermion and ghost multiplicities.
The EFT subtraction retains soft diagrams and tree Wilson insertions
generated by heavy elimination. The hard-region formulation of this
subtraction is discussed in Ref.~\cite{henning}.
In the convention of Eq.~\eqref{eq:gaussian},
\begin{equation}
 \delta K=\Delta\Gamma_{\ellight\ellight}^{(2)},\quad
 \delta J=-\Delta\Gamma_{\ellight I}^{(2)},\quad
 \delta C=-\Delta\Gamma_{II}^{(2)}.
\end{equation}
For example, the logarithm obeys
\begin{equation}
 \partial_x\partial_y\frac12\STr\log\mathcal M
 =\frac12\STr\left(\mathcal M^{-1}\mathcal M_{xy}
 -\mathcal M^{-1}\mathcal M_x\mathcal M^{-1}\mathcal M_y\right).
\end{equation}
Thus retaining only a mass Hessian is insufficient: derivatives of the
source-dependent fluctuation operator and its EFT counterpart are
essential. A one-loop shift of the stationary eliminated field does not
add a new first-order term to the stationary tree functional, but the
declared vacuum counterterms and tadpole prescription must still be used.

Our finite set of scalar-exchange current probes is not a complete local
operator basis. Full matching additionally needs appropriate Lorentz,
gauge and family structures, direct vector exchange, box contributions,
composite-current renormalization, and operator mixing. A vanishing
$C$ in a selected probe sector does not establish their absence.

\section{Actual upper heavy-vector finite kernels}
\subsection{Conventions and dimensionally finite terms}
Let $R_A$ be real antisymmetric scalar generators and $T_A$ Hermitian
fermion generators, with $T_A=iR_{F,A}$. The original Weyl convention is
\begin{equation}
 \mathcal L_Y=-\frac12\psi^TY^a\psi\,s_a+\mathrm{h.c.},\qquad
 (Y^a)^T=Y^a,\quad Z_\psi=1+k_\psi,\quad Z_s=1+k_s.
\end{equation}
Write $L_0=16\pi^2$ and
\begin{align}
 B_0(0;x,y)&=\Delta_\epsilon+b(x,y;\mu),\\
 b(x,y;\mu)&=-\int_0^1\log\frac{(1-t)x+ty}{\mu^2}\,\dd t\nonumber\\
 &=1-\frac{x\log(x/\mu^2)-y\log(y/\mu^2)}{x-y}.
\end{align}
In particular $b(x,0)=1-\log(x/\mu^2)$ and $b(x,x)=-\log(x/\mu^2)$.
The integral form controls both degenerate and one-zero limits.

The heavy Landau vector propagator is transverse. At zero external
momentum the fermion triangle gives $d-1$, whereas the scalar momentum
coefficient is $4(d-1)/d$. Therefore
\begin{align}
 [(d-1)B_0]_{\MS}&=3b-2,\\
 \left[\frac{4(d-1)}{d}B_0\right]_{\MS}&=3b-\frac12.
\end{align}
These constants would be missed by setting $d=4$ before multiplying the
UV pole. They are not adjustable matching nuisances.

For the massless fermion self-energy the numerator is
\begin{align}
 &\gamma^\mu(\not p-\not k)\gamma^\nu
 \left(g_{\mu\nu}-\frac{k_\mu k_\nu}{k^2}\right)\nonumber\\
 &\qquad=(3-d)\not p+(d-1)\not k
            -2\frac{k\cdot p}{k^2}\not k.
\end{align}
Expand the massless denominator to first order in $p$ and average
$k_\mu k_\nu\mapsto g_{\mu\nu}k^2/d$. The coefficient of $\not p$ becomes
\begin{equation}
 a(d)=5-d-\frac4d=\frac32\epsilon+O(\epsilon^2).
\end{equation}
Since the inverse propagator is $\not p-\Sigma$, its finite kinetic
coefficient in Landau gauge is
\begin{equation}\label{eq:kpsi}
 k_\psi^{V}=-\frac{3g_U^2}{2L_0}\sum_A T_A^2.
\end{equation}
The absence of a Landau logarithm in this particular kinetic graph does
not imply that the finite contribution vanishes.

\subsection{Independent Abelian mass check}
A neutral scalar mass insertion for a charged Dirac fermion provides a
check independent of the P54 group sums. Factoring out $g^2q^2/L_0$ and
putting $l=\log(M^2/\mu^2)$, the mass and slash self-energies in $R_\xi$
give
\begin{equation}
 \Sigma_m/m=1-3l+u_\xi,\qquad
 \Sigma_{\not p}=\frac32-u_\xi,\qquad
 u_\xi=\xi(1-l-\log\xi),\quad u_0=0.
\end{equation}
Thus canonical mass matching is gauge independent in this limit:
\begin{equation}
 \frac{\delta m}{m}=\frac{g^2q^2}{L_0}\left(\frac52-3l\right).
\end{equation}
The same constant/logarithm structure occurs in the neutral-vector
part of the light-fermion matching formula in Ref.~\cite{martin}.
That comparison does not import its Standard Model charges, vacuum
convention, or full threshold into P54. The executable tests also check
the $d$-dependent Laurent limits directly. RG logarithms alone cannot
verify these finite constants.

\subsection{Vertex and full active scalar metric for this subset}
Let $A$ be the $328\times248$ active embedding, $H$ the orthonormal
$328\times55$ upper heavy eigenbasis, and $m_j^2$ its positive spectrum.
The finite active vertex is
\begin{equation}\label{eq:vertex}
 \delta V^a=-\frac{g_U^2}{L_0}
 \sum_A T_A^T Y^aT_A\,[3b(M_V^2,0;\mu)-2].
\end{equation}
To avoid confusing the active embedding with a gauge index, set
$v_{A,a}=R_A A_{\cdot a}$ and define
\begin{align}
 w_0&=3b(M_V^2,0;\mu)-\tfrac12,\\
 W&=w_0(1-HH^T)+H\diag_j[3b(M_V^2,m_j^2;\mu)-\tfrac12]H^T.
\end{align}
Then the complete vector--scalar state sum for the active metric, at the
stated retained-mass expansion order, is
\begin{equation}\label{eq:ks}
 (k_s^V)_{ab}=-\frac{g_U^2}{L_0}\sum_A v_{A,a}^TWv_{A,b}.
\end{equation}
The complement is included, not deleted: it contains all retained and
zero-mode scalar directions, expanded at zero mass in the hard integral.

At this upper background, eaten directions lie in $54_H$, while active
scalars lie in $10_H+126_H$. Consequently their linear $VV s_a$ mass
vertices vanish, as do their transitions to eaten scalars through $R_A$.
Goldstone Yukawa tensors vanish because $54_H$ has no Yukawa coupling.
Mixed scalar--fermion--vector triangles at zero external momentum have a
scalar vertex proportional to vector momentum and vanish against the
Landau projector. Ghosts do not couple to fermions; active linear ghost
mass vertices vanish here. Seagulls can affect the \emph{potential}, but
not the $p^2$ metric coefficient calculated in Eq.~\eqref{eq:ks}. None of
these statements removes the separate scalar-cubic graphs or vacuum
counterterm matching.

\section{Canonical matching, Ward checks, and generated tensors}
For real scalar coordinates the first-order canonical Yukawa tensor is
\begin{equation}\label{eq:canonical}
 Y_c^a=Y^a+\delta V^a-\frac12\left[
 k_\psi^T Y^a+Y^a k_\psi+\sum_b(k_s)_{ba}Y^b\right].
\end{equation}
This external-leg structure is also obtained from finite field maps in
Ref.~\cite{patel}; its SU(5) group factors are not used here.
The real metric is essential: a complex scalar metric can include both
ordinary and conjugate-bidoublet structures after PS breaking.

The actual Clifford and scalar representations give
\begin{equation}
 C_F^{\rm coset}=3I_{16},\qquad
 C_s^{\rm coset}|_{H,F,L,R}=(3,7,6,6).
\end{equation}
Squaring the Yukawa intertwiner identity gives
\begin{equation}
 \sum_b(C_s)_{ab}Y^b=C_F^TY^a+Y^aC_F
                 +2\sum_A T_A^TY^aT_A.\label{eq:ward}
\end{equation}
Since $\dd b/\dd\log\mu=2$, Eqs.~\eqref{eq:kpsi}--\eqref{eq:canonical}
imply the independent scale check
\begin{equation}\label{eq:scale}
 \frac{\dd\delta Y_c^a}{\dd\log\mu}
 =\frac{3g_U^2}{L_0}(C_F^TY^a+Y^aC_F)
 =\left.\beta_{\rm PS}^a-\beta_{\rm Spin(10)}^a\right|_{\rm removed\ vectors}.
\end{equation}
In the auxiliary limit in which every scalar in the vector bubble is
treated at zero mass, a parent of coset Casimir $c_s$ gives
\begin{equation}
 \delta Y_c=\frac{g_U^2}{L_0}
 \left[\frac{15}{2}+\frac34 c_s
                 -9\log(M_V^2/\mu^2)\right]Y.
\end{equation}
Actual hard scalar mass splittings are retained instead of imposing this
auxiliary formula on the numerical data.

At the benchmark we obtain
\begin{align}
 k_\psi^V&=-0.00901979574931\,I_{16},\\
 \operatorname{spec}(k_s^V)&\subset
 [-0.0300659858310,-0.000821449346042].
\end{align}
Both corrected kinetic metrics $I+k$ remain positive for this subset.
The canonical finite coefficients, in the \emph{previously fixed}
six-invariant PS basis, are
\begin{center}
\begin{tabular}{llll}
\toprule
UV spurion & PS parent & Direct increment & Mirror increment\\
\midrule
$h_{\rm raw}$ & $H$ & $0.01771675106$ & $+0.0001280086893\,i$\\
$f_{\rm raw}$ & $F$ & $0.08885260368$ & $-0.0003536675766\,i$\\
$f_{\rm raw}$ & $L$ & $0.06803155989$ & none\\
$f_{\rm raw}$ & $R$ & $0.06803155989$ & none\\
\bottomrule
\end{tabular}
\end{center}
Each entry multiplies its row's UV family matrix. Other direct cross
entries vanish for this gauge contribution. Mirror phases are basis
conventions, not additional CP parameters. These linear coefficients
are distinct from the earlier cubic-spurion pure-Yukawa thresholds.
Both live at the same upper PS site, but their sum is still missing the
other upper matching pieces. The six-invariant residual is below
$4\times10^{-15}$; projecting onto only four tensors discards real terms.

\section{Numerical source-preserving transport on this background}
The physical basis is $(H,F,L,R,\mathrm{PQ},\mathrm{heavy})$, of dimension
$248+1+55=304$. At complex Euclidean $z=p_E^2$ use
\begin{equation}
 K_0(z)=B^TH_{\rm scalar}B+zI_{304},\qquad
 \delta K(z)=z\diag(k_s^V,0_{56}).
\end{equation}
The zero entries specify the \emph{chosen additive graph contribution},
not the unknown complete self-energy. The source increment consists of
the computed active vertex and the universal fermion external legs,
including on heavy scalar sources:
\begin{equation}
 \delta Y_{\rm source}^{A}=-\frac12(k_\psi^TY^A+Y^Ak_\psi)
                    +\sum_a A_{Aa}\,\delta V^a.
\end{equation}
It intentionally excludes scalar external-leg factors, since those are
already carried by $K$. Before scalar elimination $\delta C=0$ for
this additive graph set only. Unknown hard contact graphs remain unknown.

The probe ledger retains the seven P-SPEC1 spinor pairs and adds
$u^c_{1}d^c_{2}$ and $u_{L,1}d_{L,2}$, where subscripts denote the first
two matched colours in the common-phase dictionary. Each pair has
$h/f$ and real/imaginary slots, giving 36 real current slots. The old
same-colour diquark contractions vanished by antisymmetric colour
selection; they could not certify absence of the upper triplet channel.
The new action-derived sources yield nonzero upper $\delta C$.

We compare direct $304\to129$ reduction with
$304\to249\to129$ and the reverse elimination order. The final chart
retains whole $H,F$ parents and PQ. Eliminating $L,R$ here is an exact
off-shell \emph{algebra test}, not permission to decouple these physically
light parents at the upper threshold. At four nonsingular momenta,
\begin{center}
\begin{tabular}{lrr}
\toprule
$z$ & $\|\delta C_{\rm upper}\|_F$ & error after dropping final $\delta C$\\
\midrule
$0.013$ & $0.2483693$ & $0.9998932$\\
$0.27$ & $0.1336376$ & $0.9611380$\\
$-0.11+0.07i$ & $0.4106111$ & $0.6078394$\\
$0.03+0.12i$ & $0.2166991$ & $0.9651812$\\
\bottomrule
\end{tabular}
\end{center}
The error is $\|F_1-F_2\|/\max(1,\|F_1\|,\|F_2\|)$ for response
\emph{increments}, not a cross-section uncertainty. The raw family
spurions are factored out. In particular, these large diagnostic
differences do not assert large physical radiative corrections.

For an infinitesimal real-field redefinition
$\phi=(I+\ell s)\phi_c$, covariance requires
\begin{equation}
 \delta K_c=\delta K+s^TK_0+K_0s,\quad
 \delta J_c=\delta J+s^TJ_0,\quad \delta C_c=\delta C.
\end{equation}
For $s=-k_s/2$ the kinetic correction is removed while the curvature
becomes $\delta M_c=\delta M-\{k_s,M_0\}/2$.
Substitution into Eq.~\eqref{eq:deltaF} cancels all $s$ terms exactly.
Including the scalar leg in $J$ while also retaining the unnormalized
$K$ double counts it; this incorrect alternative is detected at all four
momenta. Dense synthetic jets additionally exercise nonzero $\delta B$
and $\delta D$, which happen to vanish in portions of the actual sparse
sub-contribution. They are labelled algebra tests, not physical inputs.

\section{The light-state and fit handoff, after matching closes}
In original complex doublet coordinates let
\begin{equation}
 D=D_0+\delta D_B+\delta D_F+\delta D_{\rm match}-\delta\xi P_{10},
 \qquad Z=I+k_H.
\end{equation}
Here $\delta D_F$ includes the full same-prescription fermionic tadpole
counterterm, and all Yukawa matrices are the actual evolved matched
ones. Neither unit-spurion tests nor free independent mirror matrices
can supply them. The constrained light problem is
\begin{equation}
 Dc=0,\qquad c^\dagger Zc=1,\qquad
 QD_cQ>0,\qquad D_c=Z^{-1/2}DZ^{-1/2}.
\end{equation}
The kinetic metric alone cannot rotate an exact original-coordinate null
ray. It changes its normalization and its canonical representative. For
a simple tree zero $D_0c_0=0$, $c_0^\dagger c_0=1$, define
$Q=1-c_0c_0^\dagger$ and $G=(QD_0Q)^{-1}$ on $\operatorname{ran}Q$.
First-order perturbation theory gives
\begin{align}
 \delta\xi&=\frac{c_0^\dagger\delta D_{\rm rest}c_0}
                   {c_0^\dagger P_{10}c_0},\\
 Q\delta c&=-GQ(\delta D_{\rm rest}-\delta\xi P_{10})c_0,\\
 \operatorname{Re}(c_0^\dagger\delta c)&=-\frac12c_0^\dagger k_Hc_0.
\end{align}
The phase fixes the remaining imaginary parallel part. If the corrected
heavy block is not positive or is nearly singular, perturbative
light-state promotion is not justified. An exact numerical retuning with
one-loop inputs does not create a two-loop calculation.

Up/Dirac-neutrino Yukawas contract with $c$; down/charged-lepton
vertices use the conjugate-copy dictionary and $c^*$. All six allowed
PS Yukawa tensors must run together, subject to their boundary origin in
$h_{\rm raw},f_{\rm raw}$. Subsequent nondegenerate seesaw thresholds,
$C_{HN}$, Weinberg/type-II coefficients and finite kinetic maps must use
that same trajectory. The fit objective can then be written
\begin{equation}
 \chi^2(\theta,\eta)=
 [O(\theta,\eta)-O_{\rm data}]^T V^{-1}
 [O(\theta,\eta)-O_{\rm data}]+\chi^2_{\rm nuisance}(\eta),
\end{equation}
subject to stationarity, one light doublet, the heavy-gap constraint,
perturbative control and the common RG/matching map. Any residual
$\eta_Y(\mu)$ must describe a declared uncertainty of this map, not
replace uncalculated matrices by independent per-observable parameters.
The present work cannot provide the missing map, so no minimization of
this physical objective is performed.

\section{Reproducibility, closed subgates, and next priority}
The executables are
\begin{quote}\small\ttfamily
route\_f/code/verify\_p54\_upper\_vector\_matching.py\\
route\_f/code/verify\_p54\_kjc\_finite\_matching.py
\end{quote}
Both accept \texttt{--cache-dir /path/to/tmp/p54\_full\_doublet\_cw}.
Their JSON/Markdown outputs contain matrices, fixed phases, source
hashes, all residuals and the missing-diagram ledger. The vector verifier
passes 38/38; the transport verifier passes 108/108. The pre-existing
complex canonical matching interface also passes 17/17 in memory.

The contribution composer rejects a different action, gauge,
subtraction, tadpole convention, background, scale, momentum, basis,
mass expansion or loop order, and rejects duplicated diagram labels.
The shared fit consumer now requires full upper/lower $(K,J,C)$, all-six
PS evolution, finite sequential seesaw, full fermionic light feedback and
retained-mass error control. Old boolean flags alone no longer open it.
These guards are bookkeeping checks, not proofs that a submitted diagram
inventory is physically complete.

Priority remains within the fixed action:
\begin{enumerate}
 \item Complete the upper scalar-cubic kinetic and source-dependent
 potential/tadpole/Wilson matching in this \emph{same} PS chart, including
 a controlled retained-mass expansion and a complete operator basis.
 \item Construct the lower covariant hard-minus-EFT finite functional and
 finite sequential seesaw map; do not add the old one-step broken result
 directly to the upper threshold. One-loop low-energy matching requires
 common amplitudes and their operator content, as exemplified by
 Ref.~\cite{dekens}.
 \item Only then iterate the full complex fermionic light state and the
 constrained two-spurion fit. Momentum-dependent pole masses and
 gauge/Nielsen checks remain a later, distinct requirement.
\end{enumerate}
No unrestricted theorem proving the phenomenological theory is claimed.

\begin{thebibliography}{9}
\bibitem{henning} B. Henning, X. Lu and H. Murayama,
\emph{One-loop Matching and Running with Covariant Derivative Expansion},
\href{https://arxiv.org/abs/1604.01019}{arXiv:1604.01019}.
\bibitem{martin} S. P. Martin,
\emph{Matching relations for decoupling in the Standard Model at two
loops and beyond}, Phys. Rev. D 99 (2019) 033007,
\href{https://arxiv.org/abs/1812.04100}{arXiv:1812.04100}, Eq. (2.20).
\bibitem{patel} \emph{Quantum corrections and the minimal Yukawa sector
of SU(5)}, \href{https://arxiv.org/html/2310.16563v2}{arXiv:2310.16563v2},
Appendices A/B; used for external-leg structure, not P54 group factors.
\bibitem{dekens} W. Dekens and P. Stoffer,
\emph{Low-energy effective field theory below the electroweak scale:
matching at one loop},
\href{https://arxiv.org/abs/1908.05295}{arXiv:1908.05295}.
\end{thebibliography}
\end{document}
